\chapter{Groepentheorie}\label{ch:b3:groups}

In het volume van bachelorjaar 2 dienden groepen vooral als
boekhoudkundig hulpmiddel: de stelling van Lagrange, cyclische
groepen, de symmetrische groep en haar signatuur. Dit hoofdstuk
maakt van de groepentheorie een \emph{methode}. De motor is het
begrip groep die \emph{werkt} op een verzameling: door \omterm{def:b3:groups:action}{banen} en
vaste punten te tellen krijgen we de \omterm{cor:b3:groups:classeq}{klassevergelijking}, de
stelling van Cauchy en de drie stellingen van Sylow --- het
lokaal-naar-globaalprincipe waarop de eindige groepentheorie
draait. Daarna leren we groepen samen te stellen (directe en
\omterm{def:b3:groups:semidirect}{semidirecte producten}) en weer uit elkaar te halen
(compositiereeksen, \omterm{def:b3:groups:derived}{oplosbare groepen}), en bewijzen we de stelling
die in \cref{ch:b3:galois} een drie eeuwen oude vraag over
veeltermvergelijkingen zal beslechten: de \omterm{thm:b3:groups:ansimple}{alternerende groep} $A_n$
is \omterm{def:b3:groups:simple}{enkelvoudig} voor $n \geq 5$.

\section{Quotiëntgroepen en de isomorfiestellingen}

Overal is $G$ een multiplicatief geschreven groep met neutraal
element $e$. Uit het volume van bachelorjaar 2 herinneren we ons:
deelgroepen, nevenklassen $gH$, de stelling van Lagrange ($\abs G =
[G:H]\,\abs H$ voor eindige $G$), de orde van een element, cyclische
groepen, en de symmetrische groep $S_n$ met haar
signatuurmorfisme $\varepsilon \colon S_n \to \{\pm1\}$.

\begin{definition}\label{def:b3:groups:normal}
Een deelgroep $N$ van $G$ heet \emph{normaal}\index{normaaldeler}
(notatie $N \trianglelefteq G$) als $gNg^{-1} = N$ voor elke $g
\in G$ --- gelijkwaardig: als linker- en rechternevenklassen
samenvallen, $gN = Ng$ voor alle $g$.
\end{definition}

\begin{theorem}[Quotiëntgroep]\label{thm:b3:groups:quotient}
Zij $N \trianglelefteq G$. De verzameling $G/N$ van nevenklassen,
voorzien van de vermenigvuldiging $(gN)(hN) = ghN$, is een goed
gedefinieerde groep, de \emph{quotiëntgroep}\index{quotiëntgroep},
en de \emph{kanonieke projectie} $\pi \colon G \to G/N$, $g \mapsto
gN$, is een surjectief morfisme met kern $N$. Omgekeerd is elke
kern van een groepsmorfisme \omterm{def:b3:groups:normal}{normaal}: de \omterm{def:b3:groups:normal}{normaaldelers} zijn precies
de kernen.
\end{theorem}

\begin{proof}
Alles draait om de goede definitie. Als $gN = g'N$ en $hN = h'N$,
schrijf dan $g' = gn$ en $h' = hm$ met $n, m \in N$. Dan is $g'h' =
gnhm = gh\,(h^{-1}nh)\,m \in ghN$, want $h^{-1}nh \in N$ wegens
normaliteit: het product van nevenklassen hangt niet van de
representanten af. Associativiteit, het neutrale element $eN = N$
en de inversen $(gN)^{-1} = g^{-1}N$ erft $G/N$ van $G$. Uiteraard
is $\pi$ een surjectief morfisme, en $\pi(g) = N \iff g \in N$.

Is $f \colon G \to H$ een morfisme en $k \in \ker f$, dan is
$f(gkg^{-1}) = f(g)f(k)f(g)^{-1} = e$: kernen zijn \omterm{def:b3:groups:normal}{normaal}.
\end{proof}

\begin{theorem}[Universele eigenschap; eerste isomorfiestelling]
\label{thm:b3:groups:firstiso}
Zij $f \colon G \to H$ een morfisme en $N \trianglelefteq G$ met
$N \subseteq \ker f$. Er is precies één morfisme $\bar f \colon G/N
\to H$ met $f = \bar f \circ \pi$. In het bijzonder, voor $N = \ker
f$:\index{isomorfiestellingen}
\[
G/\ker f \;\xrightarrow{\;\sim\;}\; \operatorname{im} f,
\qquad gN \mapsto f(g).
\]
\end{theorem}

\begin{proof}
Eenduidigheid: $\bar f(gN)$ moet $f(g)$ zijn. Bestaan: als $gN =
g'N$, dan is $g^{-1}g' \in N \subseteq \ker f$, dus $f(g) = f(g')$
en is $\bar f(gN) = f(g)$ goed gedefinieerd; het is een morfisme
omdat $f$ er een is. Voor $N = \ker f$ is $\bar f$ injectief, want
$\bar f(gN) = e$ betekent $g \in \ker f$, dat wil zeggen $gN = N$;
en het beeld van $\bar f$ is dat van $f$.
\end{proof}

\begin{theorem}[Tweede en derde isomorfiestelling]
\label{thm:b3:groups:secondiso}
Zij $H \leq G$ en $N \trianglelefteq G$.
\begin{enumerate}
  \item $HN = \{hn : h \in H,\, n \in N\}$ is een deelgroep, $N
        \trianglelefteq HN$, $H \cap N \trianglelefteq H$, en
        \[
        H/(H \cap N) \;\cong\; HN/N .
        \]
  \item Is bovendien $N \subseteq K \trianglelefteq G$, dan is $K/N
        \trianglelefteq G/N$ en $(G/N)\big/(K/N) \cong G/K$.
\end{enumerate}
\end{theorem}

\begin{proof}
(1) $HN$ is een deelgroep: $(hn)(h'n') = hh'\,(h'^{-1}nh')n' \in
HN$ en $(hn)^{-1} = h^{-1}(hn^{-1}h^{-1}) \in HN$, waarbij we de
normaliteit van $N$ gebruiken. Stel nu $H \hookrightarrow HN
\xrightarrow{\pi} HN/N$ samen: dit morfisme is surjectief ($hnN =
hN$) met kern $\{h \in H : h \in N\} = H \cap N$; pas
\cref{thm:b3:groups:firstiso} toe.

(2) De projectie $G/N \to G/K$, $gN \mapsto gK$, is goed
gedefinieerd ($N \subseteq K$), surjectief, met kern $K/N$; pas
opnieuw \cref{thm:b3:groups:firstiso} toe.
\end{proof}

\begin{theorem}[Correspondentiestelling]\label{thm:b3:groups:correspondence}
Zij $N \trianglelefteq G$. De afbeelding $H \mapsto H/N$ is een
bijectie tussen de deelgroepen van $G$ die $N$ bevatten en de
deelgroepen van $G/N$; ze bewaart inclusies, indices en
normaliteit, in beide richtingen.
\end{theorem}

\begin{proof}
De inverse is $\bar H \mapsto \pi^{-1}(\bar H)$. Beide
afbeeldingen sturen deelgroepen naar deelgroepen en zijn elkaars
inverse: $\pi^{-1}(H/N) = HN = H$ omdat $N \subseteq H$, en
$\pi(\pi^{-1}(\bar H)) = \bar H$ omdat $\pi$ surjectief is. Dat
inclusies bewaard blijven is duidelijk; $[G:H] = [G/N : H/N]$ omdat
$gH \mapsto (gN)(H/N)$ een goed gedefinieerde bijectie tussen de
nevenklassenverzamelingen is; en $gHg^{-1} = H$ voor alle $g$ geldt
dan en slechts dan als $(gN)(H/N)(gN)^{-1} = H/N$ voor alle $gN$,
opnieuw wegens de surjectiviteit van $\pi$.
\end{proof}

\begin{example}\label{ex:b3:groups:firstiso}
$\varepsilon \colon S_n \to \{\pm 1\}$ levert $S_n/A_n \cong
\{\pm1\}$; $\det \colon GL_n(K) \to K^\times$ levert
$GL_n(K)/SL_n(K) \cong K^\times$; $t \mapsto \eu^{2\iu\pi t}$
levert $\R/\Z \cong \mathbb U$, de cirkelgroep. Zo worden
quotiënten in de praktijk \emph{berekend}: zoek een surjectie met
de juiste kern.
\end{example}

\begin{method}\label{met:b3:groups:normal}
Om $N \trianglelefteq G$ te bewijzen, in afnemende volgorde van
elegantie: schrijf $N$ als de kern van een morfisme op $G$;
controleer $gNg^{-1} \subseteq N$ voor alle $g$ (dat volstaat: pas
het toe op $g^{-1}$ en conjugeer om de omgekeerde inclusie te
krijgen); ga na dat $N$ een vereniging van \omterm{ex:b3:groups:actions}{conjugatieklassen} is; of
merk op dat $[G:N] = 2$ (dan is $gN = Ng$ afgedwongen ---
\cref{exo:b3:groups:1}).
\end{method}

\section{Groepswerkingen}

\begin{definition}\label{def:b3:groups:action}
Een \emph{werking}\index{groepswerking} van $G$ op een verzameling
$X$ is een morfisme $\varphi \colon G \to \mathfrak{S}(X)$ naar de
groep van bijecties van $X$; men schrijft $g \cdot x$ voor
$\varphi(g)(x)$. Gelijkwaardig: een afbeelding $G \times X \to X$
met $e \cdot x = x$ en $g \cdot (h \cdot x) = (gh) \cdot x$. De
\emph{baan}\index{baan} van $x$ is $\mathcal O_x = \{g \cdot x : g
\in G\}$, de \emph{stabilisator}\index{stabilisator} is de
deelgroep $G_x = \{g : g\cdot x = x\}$, en $X^G = \{x : \forall
g,\ g \cdot x = x\}$ is de verzameling \emph{vaste punten}. De
werking heet \emph{transitief} als er precies één baan is,
\emph{trouw} als $\varphi$ injectief is, \emph{vrij} als alle
stabilisatoren triviaal zijn.
\end{definition}

\begin{example}\label{ex:b3:groups:actions}
Vijf \omterm{def:b3:groups:action}{werkingen} dragen de hele eindige groepentheorie:
\begin{enumerate}
  \item $G$ op zichzelf door \emph{linkstranslatie} $g \cdot x =
        gx$: vrij en transitief.
  \item $G$ op zichzelf door \emph{conjugatie} $g \cdot x =
        gxg^{-1}$: de \omterm{def:b3:groups:action}{banen} zijn de
        \emph{conjugatieklassen}\index{conjugatieklasse}, de
        \omterm{def:b3:groups:action}{stabilisatoren} de \emph{centralisatoren}
        \index{centralisator} $Z_G(x) = \{g : gx = xg\}$, de vaste
        punten het \emph{centrum}\index{centrum} $Z(G)$.
  \item $G$ op de nevenklassenverzameling $G/H$ door $g \cdot xH =
        gxH$: transitief, en de \omterm{def:b3:groups:action}{stabilisator} van de nevenklasse $H$
        is $H$ zelf. Elke transitieve \omterm{def:b3:groups:action}{werking} is van deze vorm
        (\cref{exo:b3:groups:8}).
  \item $G$ op de verzameling van haar deelgroepen door conjugatie:
        de \omterm{def:b3:groups:action}{stabilisator} van $H$ is de
        \emph{normalisator}\index{normalisator} $N_G(H) = \{g :
        gHg^{-1} = H\}$, de grootste deelgroep van $G$ waarin $H$
        \omterm{def:b3:groups:normal}{normaal} is.
  \item $S_n$ op $\intint{1}{n}$: de moeder van alle voorbeelden.
\end{enumerate}
\end{example}

\begin{theorem}[Baan en stabilisator]\label{thm:b3:groups:orbitstab}
De afbeelding $gG_x \mapsto g \cdot x$ is een goed gedefinieerde
bijectie $G/G_x \to \mathcal O_x$. In het bijzonder geldt voor
eindige $G$:\index{baan-stabilisatorstelling}
\[
\abs{\mathcal O_x} = [G : G_x] \quad\text{deelt } \abs G,
\]
en, omdat de \omterm{def:b3:groups:action}{banen} $X$ opdelen (het zijn de klassen van de
equivalentie $x \sim y \iff y \in \mathcal O_x$),
\[
\abs X = \sum_{i} \,[G : G_{x_i}]
\qquad (x_i\colon \text{één punt per baan}).
\]
\end{theorem}

\begin{proof}
Goed gedefinieerd en injectief: $gG_x = hG_x \iff h^{-1}g \in G_x
\iff h^{-1}g \cdot x = x \iff g \cdot x = h \cdot x$; lees de keten
in beide richtingen. Surjectiviteit is de definitie van de \omterm{def:b3:groups:action}{baan}. De
tellingen volgen uit de stelling van Lagrange en uit de opdeling
van $X$ in \omterm{def:b3:groups:action}{banen}.
\end{proof}

\begin{corollary}[Klassevergelijking]\label{cor:b3:groups:classeq}
Voor een eindige groep $G$ geldt, met één representant $x_i$ per
conjugatieklasse met meer dan één element:\index{klassevergelijking}
\[
\abs G = \abs{Z(G)} + \sum_i \,[G : Z_G(x_i)],
\qquad\text{waarbij } [G : Z_G(x_i)] > 1 \text{ en } \abs G \text{
deelt.}
\]
\end{corollary}

\begin{proof}
Pas \cref{thm:b3:groups:orbitstab} toe op de conjugatiewerking: de
\omterm{def:b3:groups:action}{banen} met één element zijn precies de elementen van $Z(G)$.
\end{proof}

\begin{theorem}[Vaste punten van $p$-groepen]\label{thm:b3:groups:pfixed}
Zij $p$ priem. Een \emph{$p$-groep}\index{p-groep@$p$-groep} is een
eindige groep waarvan de orde een macht van $p$ is. Werkt een
$p$-groep $G$ op een eindige verzameling $X$, dan is
\[
\abs{X^G} \equiv \abs X \pmod p .
\]
Gevolgen: een niet-triviale $p$-groep heeft een niet-triviaal
\omterm{ex:b3:groups:actions}{centrum}, en elke groep van orde $p^2$ is abels.
\end{theorem}

\begin{proof}
Elke \omterm{def:b3:groups:action}{baan} heeft kardinaliteit $[G:G_x]$, een macht van $p$; die
macht is $1$ precies in de vaste punten en anders deelbaar door
$p$. Sommeren over de \omterm{def:b3:groups:action}{banen} geeft de congruentie. Voor het \omterm{ex:b3:groups:actions}{centrum}:
bij de conjugatiewerking van $G$ op zichzelf is $X^G = Z(G)$, dus
$\abs{Z(G)} \equiv \abs G \equiv 0 \pmod p$, en omdat $e \in Z(G)$
volgt $\abs{Z(G)} \geq p$. Orde $p^2$: als $Z(G) \neq G$, dan is
$\abs{Z(G)} = p$ en is $G/Z(G)$ cyclisch van orde $p$, wat $G$
abels maakt (\cref{exo:b3:groups:2}) --- tegenspraak.
\end{proof}

\begin{theorem}[Cauchy]\label{thm:b3:groups:cauchy}
Deelt een priemgetal $p$ de orde $\abs G$, dan bevat $G$ een
element van orde $p$.\index{stelling van Cauchy}
\end{theorem}

\begin{proof}[Bewijs (McKay)]
Zij $X = \{(g_1, \dots, g_p) \in G^p : g_1 g_2 \cdots g_p = e\}$.
Vrij kiezen van $g_1, \dots, g_{p-1}$ legt $g_p$ vast: $\abs X =
\abs G^{p-1}$, deelbaar door $p$. De cyclische groep $\Z/p\Z$ werkt
op $X$ door cyclische verschuiving $(g_1, \dots, g_p) \mapsto (g_2,
\dots, g_p, g_1)$ --- dit laat $X$ invariant, want $g_2 \cdots g_p
g_1 = g_1^{-1}(g_1 \cdots g_p)g_1 = e$. Volgens
\cref{thm:b3:groups:pfixed} is $\abs{X^{\Z/p\Z}} \equiv \abs X
\equiv 0 \pmod p$. De vaste punten zijn de constante tupels $(g,
\dots, g)$ met $g^p = e$; het tupel $(e, \dots, e)$ is er één van,
dus zijn er minstens $p$, en bijgevolg is er minstens één $g \neq
e$ met $g^p = e$: de orde daarvan is precies $p$.
\end{proof}

\begin{theorem}[Cayley]\label{thm:b3:groups:cayley}
Elke groep van orde $n$ laat zich inbedden in $S_n$.\index{stelling
van Cayley}
\end{theorem}

\begin{proof}
De linkstranslatie $\varphi \colon G \to \mathfrak S(G) \cong S_n$
is een morfisme; uit $\varphi(g) = \mathrm{id}$ volgt $g = ge = e$:
de \omterm{def:b3:groups:action}{werking} is trouw.
\end{proof}

\begin{method}\label{met:b3:groups:counting}
Vaste punten tellen is de universele openingszet van de eindige
groepentheorie. Om te bewijzen dat iets \emph{bestaat} (een
centraal element, een element van orde $p$, een \omterm{def:b3:groups:normal}{normaaldeler}, een
vast punt), laat men een goed gekozen groep werken op een goed
gekozen eindige verzameling en vergelijkt men vervolgens
$\abs{X^G}$ met $\abs X$ modulo $p$, of laat men de baangroottes de
groepsorde delen. De bewijzen van
\cref{thm:b3:groups:pfixed,thm:b3:groups:cauchy} en van alle drie
de stellingen van Sylow hieronder zijn vijf variaties op dat ene
idee.
\end{method}

\section{De stellingen van Sylow}

De stelling van Lagrange zegt dat de orde van een deelgroep $\abs
G$ deelt; het omgekeerde faalt ($A_4$, van orde $12$, heeft geen
deelgroep van orde $6$ --- \cref{exo:b3:groups:1}). De stellingen
van Sylow redden dat omgekeerde voor priemmachten, en hun telclause
is het scherpste algemene gereedschap dat we hebben om
\omterm{def:b3:groups:normal}{normaaldelers} te produceren.

\begin{definition}\label{def:b3:groups:sylow}
Schrijf $\abs G = p^a m$ met $p \nmid m$. Een
\emph{Sylow-$p$-deelgroep}\index{Sylow-deelgroep} van $G$ is een
deelgroep van orde $p^a$ --- een $p$-deelgroep van de grootst
denkbare orde. Het aantal Sylow-$p$-deelgroepen van $G$ noteren we
$n_p$.
\end{definition}

\begin{lemma}\label{lem:b3:groups:binom}
Is $\abs G = p^a m$ met $p \nmid m$, dan geldt $\dbinom{p^a
m}{p^a} \equiv m \pmod p$.
\end{lemma}

\begin{proof}
In $\mathbb F_p[X]$ geldt de droom van de eerstejaars $(1+X)^p = 1
+ X^p$ (de coëfficiënten $\binom pk$ met $0<k<p$ zijn deelbaar door
$p$: $p$ deelt de teller van $\frac{p!}{k!(p-k)!}$ maar niet de
noemer), en herhaald toepassen geeft $(1+X)^{p^a} = 1 + X^{p^a}$,
waaruit
\[
(1+X)^{p^a m} = \bigl(1 + X^{p^a}\bigr)^m
= \sum_{k=0}^{m} \binom mk X^{k p^a} \quad\text{in } \mathbb
F_p[X].
\]
Vergelijk de coëfficiënt van $X^{p^a}$: links $\binom{p^a m}{p^a}
\bmod p$, rechts $\binom m1 = m$.
\end{proof}

\begin{theorem}[Sylow I: bestaan]\label{thm:b3:groups:sylow1}
Voor elk priemgetal $p$ bestaan er Sylow-$p$-deelgroepen van $G$.
\end{theorem}

\begin{proof}[Bewijs (Wielandt)]
Zij $\Omega$ de verzameling \emph{deelverzamelingen} van $G$ met
kardinaliteit $p^a$; $G$ werkt op $\Omega$ door linkstranslatie $g
\cdot S = gS$. Wegens \cref{lem:b3:groups:binom} is $\abs\Omega =
\binom{p^a m}{p^a} \equiv m \not\equiv 0 \pmod p$, dus is er een
\omterm{def:b3:groups:action}{baan} $\mathcal O_S$ waarvan de grootte priem is met $p$ (deelde $p$
elke baangrootte, dan deelde $p$ ook $\abs\Omega$). Zij $H = G_S$
de \omterm{def:b3:groups:action}{stabilisator} van zo'n $S$. Omdat $[G : H] = \abs{\mathcal O_S}$
priem is met $p$ en $p^a \mid \abs G = [G:H]\,\abs H$, volgt $p^a
\mid \abs H$. Omgekeerd, kies $s \in S$: de afbeelding $H \to S$,
$h \mapsto hs$, is injectief en landt in $S$ omdat $hS = S$; dus
$\abs H \leq \abs S = p^a$. Bijgevolg is $\abs H = p^a$.
\end{proof}

\begin{theorem}[Sylow II: dominantie en conjugatie]\label{thm:b3:groups:sylow2}
Zij $P$ een Sylow-$p$-deelgroep en $Q$ een willekeurige
$p$-deelgroep van $G$. Dan is $Q \subseteq gPg^{-1}$ voor zekere $g
\in G$. In het bijzonder zijn alle Sylow-$p$-deelgroepen
geconjugeerd, en $P \trianglelefteq G \iff n_p = 1$.
\end{theorem}

\begin{proof}
Laat $Q$ werken op de nevenklassenverzameling $X = G/P$, van
kardinaliteit $m \not\equiv 0 \pmod p$. Volgens
\cref{thm:b3:groups:pfixed}, toegepast op de $p$-groep $Q$, is
$\abs{X^Q} \equiv m \not\equiv 0 \pmod p$: er is een vaste
nevenklasse $gP$, dat wil zeggen $QgP = gP$, oftewel $g^{-1}Qg
\subseteq P$. Is $Q$ zelf een \omterm{def:b3:groups:sylow}{Sylow-deelgroep}, dan maakt de
gelijkheid van de ordes van $Q \subseteq gPg^{-1}$ een gelijkheid.
Ten slotte is $P \trianglelefteq G$ precies wanneer haar
conjugaten $\{gPg^{-1}\}$ --- die volgens het bovenstaande
\emph{alle} Sylow-$p$-deelgroepen zijn --- tot $\{P\}$ herleiden.
\end{proof}

\begin{theorem}[Sylow III: telling]\label{thm:b3:groups:sylow3}
Er geldt $n_p \equiv 1 \pmod p$ en $n_p = [G : N_G(P)]$, en dat
laatste deelt $m$.
\end{theorem}

\begin{proof}
Zij $\mathrm{Syl}_p$ de verzameling Sylow-$p$-deelgroepen; $G$
werkt daarop transitief door conjugatie
(\cref{thm:b3:groups:sylow2}), met als \omterm{def:b3:groups:action}{stabilisator} van $P$ de
\omterm{ex:b3:groups:actions}{normalisator} $N_G(P) \supseteq P$: dus $n_p = [G : N_G(P)]$, en $m
= [G:P] = [G:N_G(P)]\,[N_G(P):P]$ laat zien dat $n_p \mid m$.

Beperk de \omterm{def:b3:groups:action}{werking} nu tot $P$ en tel de vaste punten. Wordt $Q \in
\mathrm{Syl}_p$ vastgehouden door $P$, dan is $P \subseteq N_G(Q)$;
zowel $P$ als $Q$ is een Sylow-$p$-deelgroep van de groep $N_G(Q)$,
dus zijn ze daarin geconjugeerd (\cref{thm:b3:groups:sylow2}
toegepast op $N_G(Q)$); maar $Q \trianglelefteq N_G(Q)$, dus is $Q$
daar haar eigen enige conjugaat: $P = Q$. Het enige vaste punt is
dus $P$ zelf, en \cref{thm:b3:groups:pfixed} geeft $n_p =
\abs{\mathrm{Syl}_p} \equiv \abs{\mathrm{Syl}_p^P} = 1 \pmod p$.
\end{proof}

\begin{method}\label{met:b3:groups:sylowmethod}
Om een groep van gegeven orde $n = p^a m$ te analyseren: maak een
lijst van de delers van $m$ die congruent zijn met $1$ modulo $p$
--- dat zijn de kandidaten voor $n_p$. Is $1$ de enige kandidaat,
dan is de Sylow-$p$-deelgroep \omterm{def:b3:groups:normal}{normaal}. Wordt $n_p > 1$ afgedwongen
maar klein gehouden, laat $G$ dan door conjugatie op
$\mathrm{Syl}_p$ werken: dat levert een morfisme $G \to S_{n_p}$
met kleine kern. En tel elementen: verschillende
Sylow-$p$-deelgroepen van \emph{priemorde} $p$ snijden elkaar
triviaal, dus dragen zij $n_p(p-1)$ elementen van orde precies $p$;
overlappende tellingen voor verschillende priemgetallen dwingen
vaak een tegenspraak af (\cref{exo:b3:groups:7}).
\end{method}

\begin{example}\label{ex:b3:groups:pq}
Zij $\abs G = pq$ met $p < q$ priem en $p \nmid q - 1$. Dan dwingen
$n_q \mid p$ en $n_q \equiv 1 \bmod q$ af dat $n_q = 1$ (want $p <
q$); en $n_p \mid q$ en $n_p \equiv 1 \bmod p$ dwingen $n_p = 1$ af
(want $q \not\equiv 1 \bmod p$). Zij $P, Q$ de twee normale
\omterm{def:b3:groups:sylow}{Sylow-deelgroepen}: $P \cap Q = \{e\}$ (onderling ondeelbare ordes),
dus $\abs{PQ} = pq$ (\cref{exo:b3:groups:4}) en $G \cong P \times Q
\cong \Z/p\Z \times \Z/q\Z \cong \Z/pq\Z$ volgens
\cref{prop:b3:groups:direct} hieronder. Elke groep van orde $15$,
$33$, $35$, \dots\ is dus cyclisch. Het uitgesloten geval $p \mid q
- 1$ levert precies één groep extra, niet-abels --- zie de
weekendopgave (\cref{pb:b3:groups:1}).
\end{example}

\begin{example}[Een volledige Sylow-telling: $S_4$]
\label{ex:b3:groups:s4sylow}
Laten we de methode uitvoeren op $G = S_4$, $\abs G = 24 =
2^3\cdot3$. \emph{Sylow $3$:} $n_3 \mid 8$ en $n_3 \equiv 1 \bmod
3$, dus $n_3 \in \{1, 4\}$; omdat $\langle(123)\rangle$ en
$\langle(124)\rangle$ verschillen, is $n_3 = 4$ --- de vier
deelgroepen $\langle(abc)\rangle$, één voor elke driedelige
deelverzameling $\{a, b, c\}$, samen goed voor de $8$ drietallige
cykels. Volgens Sylow~II zijn ze geconjugeerd, en het
conjugatiemorfisme $S_4 \to S_{\mathrm{Syl}_3} \cong S_4$ is hier
een isomorfisme (de kern ligt in $N = N_G(\langle(123)\rangle)$ van
orde $24/4 = 6$, en een \omterm{def:b3:groups:normal}{normaaldeler} van $S_4$ binnen een $N$ van
$S_3$-type moet triviaal zijn: ze zou bestaan uit even permutaties
die de vier \omterm{def:b3:groups:sylow}{Sylow-deelgroepen} vasthouden, en dat doet alleen $e$).
\emph{Sylow $2$:} $n_2 \mid 3$ en $n_2 \equiv 1 \bmod 2$, dus $n_2
\in \{1, 3\}$. De deelgroep $D = \langle(1234), (13)\rangle$ heeft
orde $8$ (een \omterm{ex:b3:groups:semidirectexamples}{diëdergroep} $D_4$: de symmetrieën van het vierkant
met hoekpunten $1, 2, 3, 4$) en is niet \omterm{def:b3:groups:normal}{normaal} ($(12)(1234)(12) =
(2134)$ brengt een andere deelgroep van $4$-cykels voort), dus $n_2
= 3$: de drie kopieën van $D_4$ komen overeen met de drie manieren
om $4$ punten tot een ``vierkant'' te paren. Let op de moraal van
deze telling: $\abs{S_4} = 24$ laat beide Sylow-families de ruimte
om niet \omterm{def:b3:groups:normal}{normaal} te zijn, en beide benutten die --- vergelijk met
orde $12$, waar de telling er één van de twee \omterm{def:b3:groups:normal}{normaal} maakt
(Deel~IV van \cref{pb:b3:groups:1}).
\end{example}

\section{Producten, direct en semidirect}

\begin{proposition}[Een direct product herkennen]\label{prop:b3:groups:direct}
Zij $H, K \trianglelefteq G$ met $H \cap K = \{e\}$ en $HK = G$.
Dan is $(h,k) \mapsto hk$ een isomorfisme $H \times K \to
G$.\index{direct product}
\end{proposition}

\begin{proof}
Voor $h \in H$ en $k \in K$ ligt de \omterm{def:b3:groups:derived}{commutator} $hkh^{-1}k^{-1}$ in
$K$ (lees hem als $(hkh^{-1})k^{-1}$ en gebruik de normaliteit van
$K$) én in $H$ (lees hem als $h(kh^{-1}k^{-1})$): hij is dus $e$,
zodat $H$ en $K$ elementsgewijs commuteren en de afbeelding een
morfisme is. Ze is surjectief omdat $HK = G$, en injectief omdat
uit $hk = e$ volgt dat $h = k^{-1} \in H \cap K = \{e\}$.
\end{proof}

Het is de normaliteit van \emph{beide} factoren die het vaakst
faalt: in $S_3 = \langle (1\,2\,3)\rangle \,\langle(1\,2)\rangle$
snijden de twee factoren elkaar triviaal en brengen ze samen $S_3$
voort, en toch is $S_3 \not\cong \Z/3\Z \times \Z/2\Z$. Het juiste
begrip wanneer slechts één factor \omterm{def:b3:groups:normal}{normaal} is:

\begin{definition}\label{def:b3:groups:semidirect}
Zij $H$ en $K$ groepen en $\varphi \colon K \to
\operatorname{Aut}(H)$ een morfisme. Het \emph{semidirecte
product}\index{semidirect product} $H \rtimes_\varphi K$ is de
verzameling $H \times K$, voorzien van
\[
(h, k)\,(h', k') = \bigl(h\,\varphi(k)(h'),\; kk'\bigr).
\]
\end{definition}

\begin{proposition}\label{prop:b3:groups:semidirect}
$H \rtimes_\varphi K$ is een groep; $H \times \{e\}$ is een
\omterm{def:b3:groups:normal}{normaaldeler} isomorf met $H$ en $\{e\} \times K$ een deelgroep
isomorf met $K$; ze snijden elkaar triviaal en brengen samen het
geheel voort. Omgekeerd, als $G = NK$ met $N \trianglelefteq G$, $K
\leq G$ en $N \cap K = \{e\}$, dan is $G \cong N \rtimes_\varphi K$
met $\varphi(k) = (n \mapsto knk^{-1})$.
\end{proposition}

\begin{proof}
Rechtstreekse verificatie: de associativiteit komt neer op
$\varphi(kk') = \varphi(k)\circ\varphi(k')$ en op het feit dat elke
$\varphi(k)$ een morfisme is; het neutrale element is $(e,e)$ en
$(h,k)^{-1} = \bigl(\varphi(k^{-1})(h^{-1}), k^{-1}\bigr)$. De
projectie $(h,k) \mapsto k$ is een morfisme op $K$ met kern $H
\times \{e\}$, die daarom \omterm{def:b3:groups:normal}{normaal} is. Voor de omkering: elke $g \in
G$ schrijft zich \emph{op precies één manier} als $nk$ met $n \in
N$ en $k \in K$ (bestaan: $G = NK$; eenduidigheid: uit $nk = n'k'$
volgt $n'^{-1}n = k'k^{-1} \in N \cap K$), en
\[
(nk)(n'k') = n\,(kn'k^{-1})\;kk'
\]
laat zien dat $nk \mapsto (n, k)$ de bewerking van $G$ overbrengt
op die van $N \rtimes_\varphi K$.
\end{proof}

\begin{example}\label{ex:b3:groups:semidirectexamples}
(a) De \emph{diëdergroep}\index{diëdergroep} $D_n$ ($n \geq 3$) van
de $2n$ symmetrieën van een regelmatige $n$-hoek: de rotaties
vormen een \omterm{def:b3:groups:normal}{normaaldeler} van index $2$, elke spiegeling brengt een
complement voort, en conjugatie van een rotatie met een spiegeling
keert haar om: $D_n \cong \Z/n\Z \rtimes_\varphi \Z/2\Z$ met
$\varphi(1) = (x \mapsto -x)$.
(b) De affiene groep van een rechte, $\{x \mapsto ax + b : a \in
K^\times,\, b \in K\} \cong K \rtimes K^\times$: de translaties
zijn \omterm{def:b3:groups:normal}{normaal}, de homothetieën vormen een complement.
(c) $S_n \cong A_n \rtimes \Z/2\Z$ (complement: een willekeurige
transpositie).
(d) De \omterm{pb:b3:groups:1}{quaternionengroep} $Q_8$ is \emph{geen} \omterm{def:b3:groups:semidirect}{semidirect product}
van echte deelgroepen: elke niet-triviale deelgroep bevat $-1$
(\cref{pb:b3:groups:1}), dus snijden geen twee echte deelgroepen
elkaar triviaal.
\end{example}

\section{Oplosbare groepen; enkelvoudigheid van \texorpdfstring{$A_n$}{An}}

\begin{definition}\label{def:b3:groups:derived}
De \emph{commutator}\index{commutator} van $x, y \in G$ is $[x,y]
= xyx^{-1}y^{-1}$; de \emph{afgeleide groep}\index{afgeleide
groep} $D(G)$ is de deelgroep voortgebracht door alle
commutatoren. De \emph{afgeleide reeks} is $D^0(G) = G$,
$D^{i+1}(G) = D(D^i(G))$, en $G$ heet
\emph{oplosbaar}\index{oplosbare groep} als $D^n(G) = \{e\}$ voor
zekere $n$.
\end{definition}

\begin{proposition}\label{prop:b3:groups:derived}
$D(G)$ is \omterm{def:b3:groups:normal}{normaal} (zelfs invariant onder elk automorfisme),
$G/D(G)$ is abels, en voor $N \trianglelefteq G$ geldt: $G/N$ abels
$\iff D(G) \subseteq N$. Bovendien is $G$ \omterm{def:b3:groups:derived}{oplosbaar} dan en slechts
dan als er een keten $G = G_0 \trianglerighteq G_1
\trianglerighteq \dots \trianglerighteq G_n = \{e\}$ bestaat met
$G_{i+1} \trianglelefteq G_i$ en elk quotiënt $G_i/G_{i+1}$ abels.
Deelgroepen en quotiënten van \omterm{def:b3:groups:derived}{oplosbare groepen} zijn \omterm{def:b3:groups:derived}{oplosbaar};
omgekeerd, zijn $N$ en $G/N$ \omterm{def:b3:groups:derived}{oplosbaar}, dan ook $G$.
\end{proposition}

\begin{proof}
Een automorfisme $\alpha$ stuurt $[x,y]$ naar $[\alpha x, \alpha
y]$: het permuteert de \omterm{def:b3:groups:derived}{commutatoren} en laat dus de door hen
voortgebrachte deelgroep invariant; conjugaties zijn
automorfismen, waaruit de normaliteit volgt. In $G/D(G)$ is $\bar
x\bar y \bar x^{-1}\bar y^{-1} = \overline{[x,y]} = \bar e$: het
quotiënt is abels. Is $G/N$ abels, dan ligt elke $[x,y]$ in $N$,
dus $D(G) \subseteq N$; is omgekeerd $D(G) \subseteq N$, dan is
$G/N$ volgens de derde isomorfiestelling een quotiënt van het
abelse $G/D(G)$, en dus abels.

Is $G$ \omterm{def:b3:groups:derived}{oplosbaar}, dan is de afgeleide reeks zo'n keten. Omgekeerd,
gegeven een keten, geldt $D^i(G) \subseteq G_i$ per inductie: uit
$G_i/G_{i+1}$ abels volgt $D(G_i) \subseteq G_{i+1}$, dus
$D^{i+1}(G) = D(D^i G) \subseteq D(G_i) \subseteq G_{i+1}$; en
bijgevolg $D^n(G) = \{e\}$.

Erfelijkheid: $D^i(H) \subseteq D^i(G)$ voor $H \leq G$ (inductie),
en $D^i(G/N) = \pi(D^i(G))$ omdat $\pi$ \omterm{def:b3:groups:derived}{commutatoren} op
\omterm{def:b3:groups:derived}{commutatoren} afbeeldt; hieruit volgen de uitspraken over
deelgroepen en quotiënten. Uitbreidingen: is $D^m(G/N) = \{e\}$,
dan is $D^m(G) \subseteq N$, en uit $D^n(N) = \{e\}$ volgt
$D^{m+n}(G) = D^n(D^m(G)) \subseteq D^n(N) = \{e\}$.
\end{proof}

\begin{example}\label{ex:b3:groups:solvableexamples}
Abelse groepen zijn \omterm{def:b3:groups:derived}{oplosbaar}. $p$-groepen zijn \omterm{def:b3:groups:derived}{oplosbaar}, met
inductie op de orde: $Z(G) \neq \{e\}$ en $G/Z(G)$ is een kleinere
$p$-groep. $S_3$ en $S_4$ zijn \omterm{def:b3:groups:derived}{oplosbaar}: $S_4 \trianglerighteq A_4
\trianglerighteq V \trianglerighteq \{e\}$, waarbij $V = \{e,
(1\,2)(3\,4), (1\,3)(2\,4), (1\,4)(2\,3)\}$ de viergroep van Klein
van de dubbele transposities is (\omterm{def:b3:groups:normal}{normaal} in $S_4$: een vereniging
van \omterm{ex:b3:groups:actions}{conjugatieklassen}), met abelse quotiënten $\Z/2\Z$, $\Z/3\Z$,
$V$. In \cref{ch:b3:galois} zal ``de algemene vergelijking van
graad $n$ is \omterm{def:b3:groups:derived}{oplosbaar} door worteltrekking'' letterlijk
\emph{betekenen}: ``$S_n$ is een \omterm{def:b3:groups:derived}{oplosbare groep}''. Vandaar het
belang van de volgende definitie.
\end{example}

\begin{definition}\label{def:b3:groups:simple}
Een groep $G \neq \{e\}$ heet
\emph{enkelvoudig}\index{enkelvoudige groep} als haar enige
\omterm{def:b3:groups:normal}{normaaldelers} $\{e\}$ en $G$ zijn. Een niet-abelse enkelvoudige
groep is niet \omterm{def:b3:groups:derived}{oplosbaar}: $D(G) \trianglelefteq G$ is niet $\{e\}$
(anders was $G$ abels), dus $D(G) = G$ en is de afgeleide reeks
constant. De abelse enkelvoudige groepen zijn precies de $\Z/p\Z$
met $p$ priem (een abelse groep is enkelvoudig dan en slechts dan
als ze geen echte niet-triviale deelgroep heeft, dan en slechts dan
als ze cyclisch van priemorde is, volgens Lagrange).
\end{definition}

\begin{lemma}\label{lem:b3:groups:threecycles}
Voor $n \geq 3$ wordt $A_n$ voortgebracht door de $3$-cykels; voor
$n \geq 5$ zijn alle $3$-cykels geconjugeerd \emph{in $A_n$}.
\end{lemma}

\begin{proof}
Een element van $A_n$ is een product van een even aantal
transposities; zet ze twee aan twee en gebruik (samenstelling van
rechts naar links)
\[
(a\,b)(c\,d) = (a\,c\,b)(a\,c\,d), \qquad
(a\,b)(b\,c) = (a\,b\,c), \qquad
(a\,b)(a\,b) = e
\]
voor respectievelijk disjuncte, overlappende en gelijke paren: elk
paar transposities is een product van $3$-cykels.

Conjugatie: $\sigma(a\,b\,c)\sigma^{-1} = (\sigma a\; \sigma b\;
\sigma c)$, dus zijn twee willekeurige $3$-cykels geconjugeerd via
een zekere $\sigma \in S_n$. Is $\sigma$ oneven, vervang haar dan
door $\sigma' = \sigma (d\,e)$, waarbij $d, e$ twee punten buiten
$\{a, b, c\}$ zijn --- die bestaan omdat $n \geq 5$; dan is
$\sigma'$ even en $\sigma'(a\,b\,c) \sigma'^{-1} =
\sigma(a\,b\,c)\sigma^{-1}$, want $(d\,e)$ commuteert met
$(a\,b\,c)$.
\end{proof}

\begin{theorem}[Enkelvoudigheid van de alternerende groep]\label{thm:b3:groups:ansimple}
$A_n$ is \omterm{def:b3:groups:simple}{enkelvoudig} voor $n \geq 5$.\index{alternerende groep}
\end{theorem}

\begin{proof}
Zij $N \trianglelefteq A_n$ met $N \neq \{e\}$. Wegens
\cref{lem:b3:groups:threecycles} volstaat het aan te tonen dat $N$
\emph{één} $3$-cykel bevat: normaliteit en de conjugatie van alle
$3$-cykels in $A_n$ leggen dan alle $3$-cykels in $N$, zodat $N =
A_n$.

Voor $\rho \in S_n$ noemen we $F(\rho) = \{x : \rho(x) \neq x\}$ de
\emph{drager} van $\rho$ en $f(\rho) = \abs{F(\rho)}$. Kies $\sigma
\in N \setminus \{e\}$ met $f(\sigma)$ \emph{minimaal}. Merk op dat
een niet-triviale even permutatie $f \geq 3$ heeft, en dat
$f(\sigma) = 4$ voor $\sigma \in A_n$ alleen mogelijk is als
$\sigma$ een dubbele transpositie is (een $4$-cykel is oneven). We
tonen aan dat $\sigma$ een $3$-cykel is.

\emph{Geval A: $\sigma$ is een product van disjuncte
transposities}, zeg $\sigma = (a\,b)(c\,d)\cdots$ met $f(\sigma)
\geq 4$. Kies $e' \notin \{a, b, c, d\}$ (mogelijk, want $n \geq
5$), stel $\tau = (c\,d\,e')$ en
\[
\sigma' = \tau\sigma\tau^{-1}\,\sigma^{-1} \in N
\qquad (\tau\sigma\tau^{-1} \in N \text{ wegens normaliteit}).
\]
Omdat $\sigma\tau^{-1}\sigma^{-1} = (\sigma c\;\sigma e'\;\sigma d)
= (d\;\sigma e'\;c)$ (met $\sigma c = d$ en $\sigma d = c$), krijgen
we $\sigma' = (c\,d\,e')(d\;\sigma e'\;c)$.

Als $\sigma e' = e'$ (wat in het bijzonder geldt wanneer $f(\sigma)
= 4$, dat wil zeggen $\sigma = (a\,b)(c\,d)$): dan is $(d\,e'\,c) =
(c\,d\,e')$ en $\sigma' = (c\,d\,e')^2 = (c\,e'\,d)$, een $3$-cykel
in $N$, met $f(\sigma') = 3 < 4 \leq f(\sigma)$ --- in strijd met
de minimaliteit.

Als $\sigma e' \neq e'$: dan is $\sigma e' \notin \{a, b, c, d,
e'\}$ ($\sigma$ verwisselt $a,b$ en $c,d$, en $e' \notin
\{a,b,c,d\}$ terwijl $\sigma$ injectief is), dus beweegt $\sigma$
de zes punten $a, b, c, d, e', \sigma e'$: $f(\sigma) \geq 6$.
Anderzijds voldoet $\sigma'$, een product van twee $3$-cykels met
dragers in $\{c, d, e', \sigma e'\}$, aan $f(\sigma') \leq 4$; en
$\sigma' \neq e$, want $\sigma'(d) = \tau\sigma\tau^{-1}(c) =
\tau\sigma(e') = \sigma e' \neq d$ ($\tau$ houdt $\sigma e' \notin
\{c,d,e'\}$ vast). Dus is $\sigma' \in N \setminus\{e\}$ met
$f(\sigma') \leq 4 < f(\sigma)$: opnieuw in strijd met de
minimaliteit.

\emph{Geval B: een cykel van $\sigma$ heeft lengte $\geq 3$}, zeg
$\sigma(a) = b$ en $\sigma(b) = c$ met $a, b, c$ verschillend. Is
$\sigma$ precies die $3$-cykel, dan zijn we klaar. Zo niet, dan is
$f(\sigma) \geq 5$ (het geval $f(\sigma) = 4$ met een cykel van
lengte $\geq 3$ is de oneven $4$-cykel, uitgesloten), zodat we $d,
e' \in F(\sigma) \setminus \{a, b, c\}$ kunnen kiezen. Stel $\tau =
(c\,d\,e')$ en $\sigma' = \tau\sigma\tau^{-1}\sigma^{-1} \in N$.
Net als hiervoor beweegt $\sigma' = (c\,d\,e')\,(\sigma c\;\sigma
e'\;\sigma d)$ alleen punten van
\[
M = \{c, d, e'\} \cup \{\sigma c, \sigma d, \sigma e'\}
\subseteq F(\sigma)
\]
(beelden van bewogen punten worden bewogen: uit $\sigma(x) \ne x$
volgt $\sigma(\sigma x) \neq \sigma x$, want $\sigma$ is
injectief). Verder is $b \notin M$: de vijf punten $a, b, c, d, e'$
zijn verschillend, dus $b \notin \{c, d, e'\}$; en $b \in \{\sigma
c, \sigma d, \sigma e'\}$ zou $a \in \{c, d, e'\}$ afdwingen (pas
$\sigma^{-1}$ toe en gebruik $\sigma a = b$), en dat is onjuist.
Bijgevolg houdt $\sigma'$ het punt $b$ vast terwijl $\sigma$ het
beweegt; en $F(\sigma') \subseteq F(\sigma)$. Ten slotte is
$\sigma' \neq e$: $\sigma^{-1}(c) = b$, $\tau^{-1}(b) = b$,
$\sigma(b) = c$ en $\tau(c) = d$, dus $\sigma'(c) = d \neq c$. Dus
$\sigma' \in N \setminus \{e\}$ met $f(\sigma') \leq f(\sigma) -
1$, in strijd met de minimaliteit.

Beide gevallen zijn onmogelijk, dus is $\sigma$ een $3$-cykel.
\end{proof}

\begin{corollary}\label{cor:b3:groups:snnotsolvable}
Voor $n \geq 5$ zijn $A_n$ en $S_n$ niet \omterm{def:b3:groups:derived}{oplosbaar}, en zijn de
enige \omterm{def:b3:groups:normal}{normaaldelers} van $S_n$ de groepen $\{e\}$, $A_n$ en $S_n$.
\end{corollary}

\begin{proof}
$A_n$ is niet-abels en \omterm{def:b3:groups:simple}{enkelvoudig}, dus niet \omterm{def:b3:groups:derived}{oplosbaar}
(\cref{def:b3:groups:simple}); en een groep die een niet-oplosbare
deelgroep bevat, is zelf niet \omterm{def:b3:groups:derived}{oplosbaar}
(\cref{prop:b3:groups:derived}). Zij $N \trianglelefteq S_n$: dan
is $N \cap A_n \trianglelefteq A_n$ gelijk aan $\{e\}$ of aan
$A_n$. Is $N \cap A_n = A_n$, dan is $A_n \subseteq N$ en dwingt de
index $N \in \{A_n, S_n\}$ af. Is $N \cap A_n = \{e\}$, dan is de
beperking tot $N$ van de projectie $S_n \to S_n/A_n \cong \Z/2\Z$
injectief, dus $\abs N \leq 2$; is $N = \{e, \sigma\}$, dan maakt
de normaliteit de conjugatieklasse van $\sigma$ gelijk aan
$\{\sigma\}$, dat wil zeggen $\sigma \in Z(S_n)$. Maar $Z(S_n) =
\{e\}$ voor $n \geq 3$: beweegt $\sigma \ne e$ het punt $a$ naar $b
\neq a$, kies dan $c \notin \{a, b\}$; dan stuurt
$(b\,c)\sigma(b\,c)^{-1}$ het punt $a$ naar $c \ne b$, zodat deze
permutatie van $\sigma$ verschilt. Bijgevolg is $N = \{e\}$.
\end{proof}

\begin{theorem}[Jordan--Hölder]\label{thm:b3:groups:jordanholder}
Elke eindige groep $G \neq \{e\}$ bezit een
\emph{compositiereeks}\index{compositiereeks}
\[
\{e\} = G_0 \trianglelefteq G_1 \trianglelefteq \cdots
\trianglelefteq G_r = G,
\qquad G_{i}/G_{i-1} \text{ enkelvoudig},
\]
en de multiverzameling van de \emph{compositiefactoren}
$G_i/G_{i-1}$ hangt, op isomorfie na, niet van de gekozen reeks af.
Een eindige groep is \omterm{def:b3:groups:derived}{oplosbaar} dan en slechts dan als al haar
compositiefactoren cyclisch van priemorde zijn.\index{stelling van
Jordan--Hölder}
\end{theorem}

\begin{proof}
\emph{Bestaan}: inductie op $\abs G$. Is $G$ \omterm{def:b3:groups:simple}{enkelvoudig}, neem dan
$\{e\} \trianglelefteq G$. Kies anders een maximale echte
\omterm{def:b3:groups:normal}{normaaldeler} $N$ (er zijn eindig veel deelgroepen); $G/N$ is
\omterm{def:b3:groups:simple}{enkelvoudig} volgens de correspondentiestelling (een echte
niet-triviale \omterm{def:b3:groups:normal}{normaaldeler} van $G/N$ zou zich optillen tot een
\omterm{def:b3:groups:normal}{normaaldeler} van $G$ strikt tussen $N$ en $G$). Voeg $N
\trianglelefteq G$ toe aan een \omterm{thm:b3:groups:jordanholder}{compositiereeks} van $N$.

\emph{Eenduidigheid}: inductie op $\abs G$, waarbij het geval $G$
\omterm{def:b3:groups:simple}{enkelvoudig} duidelijk is. Neem twee compositiereeksen, met
voorlaatste termen $M \trianglelefteq G$ en $N \trianglelefteq G$
(dus $G/M$ en $G/N$ zijn \omterm{def:b3:groups:simple}{enkelvoudig}). Is $M = N$, dan besluiten we
met de inductiehypothese toegepast op $M$. Anders is $MN$, \omterm{def:b3:groups:normal}{normaal}
in $G$ en $M$ strikt bevattend, gelijk aan $G$ ($M$ is maximaal
\omterm{def:b3:groups:normal}{normaal}: elke normale $L$ met $M \subsetneq L \subsetneq G$ zou
afgebeeld worden op een echte niet-triviale \omterm{def:b3:groups:normal}{normaaldeler} van het
enkelvoudige $G/M$). De tweede isomorfiestelling geeft
\[
G/M = MN/M \cong N/(M \cap N),
\qquad
G/N = MN/N \cong M/(M \cap N).
\]
Stel $K = M \cap N$ ($\trianglelefteq G$) en kies een
\omterm{thm:b3:groups:jordanholder}{compositiereeks} van $K$. Dan draagt $M$ twee compositiereeksen: de
oorspronkelijke, en de reeks van $K$ gevolgd door $K
\trianglelefteq M$ (het quotiënt $M/K \cong G/N$ is \omterm{def:b3:groups:simple}{enkelvoudig}).
Per inductie (toegepast op $M$) zijn de factoren van de
oorspronkelijke reeks van $M$ dus $\{\text{factoren van } K\} \cup
\{G/N\}$; en evenzo voor $N$. Bijgevolg hebben beide reeksen van
$G$ de factoren
\[
\{\text{factoren van } K\} \;\cup\; \{\,G/N,\; G/M\,\},
\]
dezelfde multiverzameling.

\omterm{def:b3:groups:derived}{Oplosbaarheid}: zijn alle factoren $\Z/p_i\Z$, dan is de reeks een
keten met abelse quotiënten, dus is $G$ \omterm{def:b3:groups:derived}{oplosbaar}
(\cref{prop:b3:groups:derived}). Omgekeerd is een compositiefactor
van een \omterm{def:b3:groups:derived}{oplosbare groep} zelf \omterm{def:b3:groups:derived}{oplosbaar} (een quotiënt van een
deelgroep) én \omterm{def:b3:groups:simple}{enkelvoudig}; en een oplosbare \omterm{def:b3:groups:simple}{enkelvoudige groep} is
abels (uit $D(G) \ne G$ volgt $D(G) = \{e\}$), dus van de vorm
$\Z/p\Z$.
\end{proof}

\begin{remark}\label{rem:b3:groups:jhmeaning}
Jordan--Hölder zegt dat elke eindige groep is opgebouwd uit
\omterm{def:b3:groups:simple}{enkelvoudige groepen}, met een welbepaalde stuklijst --- een
rekenkunde van groepen waarin de \omterm{def:b3:groups:simple}{enkelvoudige groepen} de
priemgetallen zijn, en waarin \emph{de manier waarop de onderdelen
aan elkaar zitten} (de uitbreidingsgegevens, zoals bij het
\omterm{def:b3:groups:semidirect}{semidirecte product}) de plaats inneemt van een gewone
vermenigvuldiging. De classificatie van de eindige \omterm{def:b3:groups:simple}{enkelvoudige
groepen} --- de cyclische $\Z/p\Z$, de alternerende $A_{n \geq 5}$,
zestien families van Lie-type en $26$ sporadische groepen --- is
een van de monumenten van de twintigste-eeuwse wiskunde; het
bewijs, verspreid over zo'n tienduizend tijdschriftpagina's, gaat
deze cursus ver te boven.
\end{remark}

\begin{omfigure}
\begin{tikzpicture}[xscale=1.55, yscale=1.4]
  \tikzset{sg/.style={draw, rounded corners=2pt, inner sep=3pt,
      fill=omDef!6, font=\small}}
  \node[sg, fill=omThm!10] (G)   at (0, 3)    {$D_4$};
  \node[sg] (K1)  at (-1.6, 2)   {$\langle r^2, s\rangle$};
  \node[sg] (R)   at (0, 2)      {$\langle r\rangle$};
  \node[sg] (K2)  at (1.6, 2)    {$\langle r^2, rs\rangle$};
  \node[sg] (S1)  at (-2.4, 1)   {$\langle s\rangle$};
  \node[sg] (S2)  at (-1.2, 1)   {$\langle r^2 s\rangle$};
  \node[sg, fill=omProp!12] (Z) at (0, 1) {$\langle r^2\rangle$};
  \node[sg] (S3)  at (1.2, 1)    {$\langle rs\rangle$};
  \node[sg] (S4)  at (2.4, 1)    {$\langle r^3 s\rangle$};
  \node[sg] (E)   at (0, 0)      {$\{e\}$};
  \draw (G) -- (K1); \draw (G) -- (R); \draw (G) -- (K2);
  \draw (K1) -- (S1); \draw (K1) -- (S2); \draw (K1) -- (Z);
  \draw (R) -- (Z);
  \draw (K2) -- (Z); \draw (K2) -- (S3); \draw (K2) -- (S4);
  \draw (S1) -- (E); \draw (S2) -- (E); \draw (Z) -- (E);
  \draw (S3) -- (E); \draw (S4) -- (E);
\end{tikzpicture}

{\small De tien deelgroepen van de \omterm{ex:b3:groups:semidirectexamples}{diëdergroep} $D_4 = \langle r,
s \mid r^4 = s^2 = e,\ srs^{-1} = r^{-1}\rangle$. De drie
deelgroepen van index $2$ (middelste rij) zijn \omterm{def:b3:groups:normal}{normaal}, net als het
\omterm{ex:b3:groups:actions}{centrum} $\langle r^2\rangle$ (gemarkeerd); de vier
spiegelingsdeelgroepen vallen uiteen in twee \omterm{ex:b3:groups:actions}{conjugatieklassen} van
twee. Ketens van onder naar boven geven compositiereeksen,
bijvoorbeeld $\{e\} \trianglelefteq \langle r^2\rangle
\trianglelefteq \langle r\rangle \trianglelefteq D_4$: factoren
$\Z/2\Z, \Z/2\Z, \Z/2\Z$ --- steeds dezelfde multiverzameling,
zoals Jordan--Hölder eist.}
\end{omfigure}

\section{Oefeningen}

\begin{exercise}[$\star$]\label{exo:b3:groups:1}
(a) Toon aan dat elke deelgroep van index $2$ \omterm{def:b3:groups:normal}{normaal} is.
(b) Toon aan dat uit $[G : H] = 2$ volgt dat $x^2 \in H$ voor elke
$x \in G$.
(c) Leid af dat $A_4$ geen deelgroep van orde $6$ heeft: het
omgekeerde van de stelling van Lagrange faalt. \emph{(Tel de
kwadraten van de $3$-cykels.)}
\end{exercise}

\begin{exercise}[$\star$]\label{exo:b3:groups:2}
Toon aan dat $G$ abels is zodra $G/Z(G)$ cyclisch is. Leid daaruit
opnieuw af dat elke groep van orde $p^2$ abels is, en geef voor elk
priemgetal $p$ een niet-abelse groep van orde $p^3$. \emph{(Denk
aan bovendriehoeksmatrices met eenheidsdiagonaal over $\mathbb
F_p$.)}
\end{exercise}

\begin{exercise}[$\star$]\label{exo:b3:groups:3}
(a) Toon aan dat $\operatorname{Aut}(\Z/n\Z) \cong
(\Z/n\Z)^\times$.
(b) Toon aan dat de inwendige automorfismen $\iota_g \colon x
\mapsto gxg^{-1}$ een \omterm{def:b3:groups:normal}{normaaldeler} $\operatorname{Inn}(G)
\trianglelefteq \operatorname{Aut}(G)$ vormen, met
$\operatorname{Inn}(G) \cong G/Z(G)$.
\end{exercise}

\begin{exercise}[$\star\star$]\label{exo:b3:groups:4}
Zij $H, K$ deelgroepen van een eindige groep $G$.
(a) Bewijs de \emph{productformule} $\abs{HK}\,\abs{H\cap K} =
\abs H\, \abs K$ door de vezels van de afbeelding $H \times K \to
HK$, $(h,k) \mapsto hk$, te tellen.
(b) Toon aan dat $HK$ een deelgroep is dan en slechts dan als $HK =
KH$ (wat automatisch geldt als een van beide \omterm{def:b3:groups:normal}{normaal} is).
(c) Toon aan dat uit $H, K \trianglelefteq G$ en $H \cap K = \{e\}$
volgt dat $hk = kh$ voor alle $h \in H$ en $k \in K$.
\end{exercise}

\begin{exercise}[$\star\star$]\label{exo:b3:groups:5}
(Teltelling van Burnside) Een eindige groep $G$ werkt op een
eindige verzameling $X$. Toon aan dat het aantal \omterm{def:b3:groups:action}{banen} het
\emph{gemiddelde aantal vaste punten} is:
\[
\#\{\text{banen}\} = \frac{1}{\abs G}\sum_{g \in G}
\abs{\operatorname{Fix}(g)},
\qquad \operatorname{Fix}(g) = \{x \in X : g \cdot x = x\},
\]
door de verzameling $\{(g,x) : g\cdot x = x\}$ op twee manieren te
tellen. Toepassing: $\Z/p\Z$ ($p$ priem) werkt door rotatie op
kettingen van $p$ kralen in $a$ beschikbare kleuren; leid daaruit
de kleine stelling van Fermat af, $a^p \equiv a \pmod p$.
\end{exercise}

\begin{exercise}[$\star$]\label{exo:b3:groups:6}
Toon met de stellingen van Sylow aan dat elke groep van orde $15$
cyclisch is en dat elke groep van orde $45$ abels is.
\end{exercise}

\begin{exercise}[$\star\star$]\label{exo:b3:groups:7}
Toon aan dat geen enkele groep van orde $30$ en geen enkele van
orde $56$ \omterm{def:b3:groups:simple}{enkelvoudig} is. \emph{(Voor $30$: tel, in het geval $n_3
\neq 1$ en $n_5 \neq 1$, de elementen van orde $3$ en $5$. Voor
$56$: tel de elementen van orde $7$.)}
\end{exercise}

\begin{exercise}[$\star\star$]\label{exo:b3:groups:8}
(a) Zij $H \leq G$ van index $n$. Toon aan dat de \omterm{def:b3:groups:action}{werking} van $G$
op $G/H$ een morfisme $G \to S_n$ oplevert waarvan de kern
$\bigcap_{g \in G} gHg^{-1}$ de grootste in $H$ bevatte
\omterm{def:b3:groups:normal}{normaaldeler} van $G$ is.
(b) Leid af: is $G$ eindig en is $p$ de \emph{kleinste} priemdeler
van $\abs G$, dan is elke deelgroep van index $p$ \omterm{def:b3:groups:normal}{normaal}.
(c) Toon aan dat elke transitieve \omterm{def:b3:groups:action}{werking} van $G$ op een
verzameling $X$ isomorf is met de \omterm{def:b3:groups:action}{werking} op een
nevenklassenverzameling: er bestaat een bijectie $X \to G/G_x$ die
met de \omterm{def:b3:groups:action}{werkingen} verwisselbaar is.
\end{exercise}

\begin{exercise}[$\star\star$]\label{exo:b3:groups:9}
(a) Toon aan dat $D(G)$ de kleinste \omterm{def:b3:groups:normal}{normaaldeler} van $G$ met abels
quotiënt is, en dat elk morfisme van $G$ naar een abelse groep op
precies één manier over de \emph{abelianisering}
$G^{\mathrm{ab}} = G/D(G)$ factoriseert.
(b) Bereken $D(S_n)$ en $S_n^{\mathrm{ab}}$ voor $n \geq 2$, en
$D(Q_8)$ en $Q_8^{\mathrm{ab}}$.
\end{exercise}

\begin{exercise}[$\star\star$]\label{exo:b3:groups:10}
Zij $G$ een $p$-groep en $H \subsetneq G$ een echte deelgroep. Toon
aan dat $H \subsetneq N_G(H)$ (``\omterm{ex:b3:groups:actions}{normalisatoren} groeien'') en leid
af dat elke maximale deelgroep van een $p$-groep \omterm{def:b3:groups:normal}{normaal} is van
index $p$. \emph{(Inductie op $\abs G$ met $Z(G) \neq \{e\}$:
behandel $Z(G) \subseteq H$ en $Z(G) \not\subseteq H$
afzonderlijk.)}
\end{exercise}

\begin{exercise}[$\star\star\star$]\label{exo:b3:groups:11}
(Enkelvoudigheid van $A_5$, met de handen)
(a) Toon aan dat de \omterm{ex:b3:groups:actions}{conjugatieklassen} van $A_5$ de kardinaliteiten
$1$, $15$, $20$, $12$, $12$ hebben. Let op het uiteenvallen van de
$S_5$-klasse van $5$-cykels: vergelijk voor een $5$-cykel $\sigma$
de \omterm{ex:b3:groups:actions}{centralisatoren} van $\sigma$ in $S_5$ en in $A_5$.
(b) Leid af dat $A_5$ \omterm{def:b3:groups:simple}{enkelvoudig} is: een \omterm{def:b3:groups:normal}{normaaldeler} is een
vereniging van \omterm{ex:b3:groups:actions}{conjugatieklassen}, bevat $e$ en heeft een
kardinaliteit die $60$ deelt.
(c) Toon aan dat een \omterm{def:b3:groups:simple}{enkelvoudige groep} van orde $60$ noodzakelijk
$n_5 = 6$ heeft.
\end{exercise}

\begin{exercise}[$\star\star$]\label{exo:b3:groups:12}
(\omterm{ex:b3:groups:actions}{Normalisatoren} van \omterm{def:b3:groups:sylow}{Sylow-deelgroepen} zijn hun eigen \omterm{ex:b3:groups:actions}{normalisator})
Zij $P$ een Sylow-$p$-deelgroep van een eindige groep $G$ en $H =
N_G(P)$.
(a) Toon aan dat $P$ de \emph{enige} Sylow-$p$-deelgroep van $H$
is.
(b) Leid af dat $N_G(H) = H$. \emph{(Voor $g \in N_G(H)$ is
$gPg^{-1}$ een Sylow-$p$-deelgroep van $H$, dus $gPg^{-1} = P$.)}
(c) Besluit dat geen enkele \omterm{ex:b3:groups:actions}{Sylow-normalisator} in een echte
\omterm{def:b3:groups:normal}{normaaldeler} van $G$ ligt, en dat een maximale deelgroep die
$N_G(P)$ bevat haar eigen \omterm{ex:b3:groups:actions}{normalisator} is.
\end{exercise}

\section{Probleem: de groepen van orde ten hoogste 15}

\begin{problem}[{Weekendopgave --- classificatie van kleine
groepen}]\label{pb:b3:groups:1}
Het doel is een volledige classificatie, met volledige bewijzen,
van de groepen van orde $\leq 15$ op isomorfie na. De ordes $1, 2,
3, 5, 7, 11, 13$ zijn afgehandeld door Lagrange (cyclisch), en de
ordes $4$ en $9$ door \cref{thm:b3:groups:pfixed} samen met de
analyse van $p^2$ hieronder: blijven over $6, 8, 10, 12, 14, 15$.

\medskip
\noindent\textbf{Deel I --- Gereedschap.}
\begin{enumerate}
  \item Toon aan dat een groep waarin elk element aan $x^2 = e$
        voldoet abels is; leid af dat zo'n eindige groep orde $2^k$
        heeft en isomorf is met $(\Z/2\Z)^k$. \emph{(Beschouw haar
        als een vectorruimte over $\mathbb F_2$.)}
  \item Toon aan dat een groep van orde $p^2$ isomorf is met
        $\Z/p^2\Z$ of met $(\Z/p\Z)^2$. Geef de abelse groepen van
        orde $8$ op isomorfie na: $\Z/8\Z$, $\Z/4\Z \times \Z/2\Z$,
        $(\Z/2\Z)^3$ --- bewijs dat die lijst volledig is en geen
        dubbels bevat, \emph{zonder} de structuurstelling van
        \cref{ch:b3:modules} (onderscheid de gevallen naar de
        maximale orde van een element).
  \item Zij $\varphi, \varphi' \colon K \to \operatorname{Aut}(H)$
        twee \omterm{def:b3:groups:action}{werkingen}. Toon aan dat $H \rtimes_{\varphi} K \cong H
        \rtimes_{\varphi'} K$ zodra $\varphi' = \varphi \circ
        \alpha$ met $\alpha \in \operatorname{Aut}(K)$.
  \item Bepaal $\operatorname{Aut}(\Z/n\Z)$ expliciet voor $n = 3,
        4, 5, 7$, en toon aan dat
        $\operatorname{Aut}\bigl((\Z/2\Z)^2 \bigr) \cong S_3$.
\end{enumerate}

\medskip
\noindent\textbf{Deel II --- Ordes $2p$ ($6$, $10$, $14$) en
$pq$.}
\begin{enumerate}[resume]
  \item Zij $\abs G = 2p$ met $p$ een oneven priemgetal. Toon aan
        dat $G$ een \omterm{def:b3:groups:normal}{normaaldeler} $N = \langle r \rangle$ van orde
        $p$ bezit en een element $s$ van orde $2$ buiten $N$.
  \item Leid af dat $G \cong \Z/p\Z \rtimes_\varphi \Z/2\Z$, waarbij
        $\varphi(1) \in \operatorname{Aut}(\Z/p\Z)$ een involutie
        is, en besluit: $G \cong \Z/2p\Z$ of $G \cong D_p$; ga na
        dat die twee niet isomorf zijn. Hiermee zijn de ordes $6$,
        $10$ en $14$ afgehandeld.
  \item Zij algemener $\abs G = pq$ met $p < q$ priem. Toon aan dat
        $G$ cyclisch is als $p \nmid q-1$ (\cref{ex:b3:groups:pq}),
        en dat er bij $p \mid q - 1$ naast $\Z/pq\Z$ \emph{precies
        één} niet-abelse groep $\Z/q\Z \rtimes \Z/p\Z$ bestaat op
        isomorfie na --- gebruik vraag 3 en het feit dat
        $\operatorname{Aut}(\Z/q\Z) \cong (\Z/q\Z)^\times$ cyclisch
        van orde $q - 1$ is, hier zonder bewijs aangenomen en
        bewezen in \cref{ch:b3:galois} (cycliciteit van $\mathbb
        F_q^\times$). Besluit voor orde $15$.
\end{enumerate}

\medskip
\noindent\textbf{Deel III --- Orde $8$.} Zij $G$ niet-abels van
orde $8$.
\begin{enumerate}[resume]
  \item Toon aan dat $G$ een element $r$ van orde $4$ bezit
        (gebruik vraag 1) en dat $N = \langle r\rangle$ \omterm{def:b3:groups:normal}{normaal} is.
  \item Zij $s \notin N$. Toon aan dat $s^2 \in N$
        (\cref{exo:b3:groups:1}(b)), dat $srs^{-1} = r^{-1}$
        (bekijk de mogelijke beelden van $r$ onder conjugatie, die
        orde $4$ moeten hebben, en sluit $srs^{-1} = r$ uit), en
        dat $s^2 \in \{e, r^2\}$ \emph{(wat gebeurt er als $s^2 =
        r$ of $r^3$? en waarom moet $s^2$ met $s$ commuteren?)}.
  \item Toon in het geval $s^2 = e$ aan dat $G \cong D_4$.
  \item Toon in het geval $s^2 = r^2$ aan dat de
        vermenigvuldigingstabel volledig vastligt; de resulterende
        groep is de
        \emph{quaternionengroep}\index{quaternionengroep} $Q_8 =
        \{\pm 1, \pm \mathrm i, \pm \mathrm j, \pm \mathrm k\}$ met
        $\mathrm i^2 = \mathrm j^2 = \mathrm k^2 = \mathrm i\,
        \mathrm j\,\mathrm k = -1$ (stel $r = \mathrm i$ en $s =
        \mathrm j$). Ga na dat $Q_8$ bestaat, bijvoorbeeld binnen
        $GL_2(\C)$ via
        \[
        \mathrm i \mapsto \begin{pmatrix} \iu & 0\\ 0 & -\iu
        \end{pmatrix},
        \qquad
        \mathrm j \mapsto \begin{pmatrix} 0 & 1\\ -1 & 0
        \end{pmatrix}.
        \]
  \item Toon aan dat elke niet-triviale deelgroep van $Q_8$ het
        element $-1$ bevat; leid af dat elke deelgroep van $Q_8$
        \omterm{def:b3:groups:normal}{normaal} is, dat $D_4 \not\cong Q_8$ (tel de elementen van
        orde $2$), en dat $Q_8$ geen \omterm{def:b3:groups:semidirect}{semidirect product} van twee
        echte deelgroepen is.
\end{enumerate}

\medskip
\noindent\textbf{Deel IV --- Orde $12$.} Zij $\abs G = 12$, $P_3
\in \mathrm{Syl}_3(G)$ en $P_2 \in \mathrm{Syl}_2(G)$.
\begin{enumerate}[resume]
  \item Toon aan dat $n_3 \in \{1, 4\}$ en $n_2 \in \{1, 3\}$, en
        dat $n_3 = 4$ afdwingt dat $n_2 = 1$ \emph{(tel de
        elementen van orde $3$)}.
  \item Stel $n_3 = 4$. De conjugatiewerking op $\mathrm{Syl}_3$
        geeft $\rho \colon G \to S_4$. Toon aan dat $\ker \rho$,
        bevat in elke $N_G(P_3)$ en dus van een orde die $3$ deelt,
        triviaal is \emph{(waarom kan hij geen orde $3$ hebben?)};
        dat het beeld, een deelgroep van orde $12$ van $S_4$,
        noodzakelijk $A_4$ is \emph{(deelgroepen van index $2$ zijn
        \omterm{def:b3:groups:normal}{normaal} en bevatten alle kwadraten ---
        \cref{exo:b3:groups:1}; tel de kwadraten in $S_4$)}; en
        besluit dat $G \cong A_4$.
  \item Stel $n_3 = 1$, zodat $G \cong \Z/3\Z \rtimes_\varphi P_2$
        met $\varphi \colon P_2 \to \operatorname{Aut}(\Z/3\Z)
        \cong \Z/2\Z$. Loop de gevallen af: voor triviale $\varphi$
        krijgen we $\Z/12\Z$ en $\Z/6\Z \times \Z/2\Z$; voor $P_2 =
        \Z/4\Z$ met surjectieve $\varphi$ de \emph{dicyclische
        groep} $\mathrm{Dic}_3 = \Z/3\Z \rtimes \Z/4\Z$; en voor
        $P_2 = (\Z/2\Z)^2$ met surjectieve $\varphi$, op de
        gelijkwaardigheid uit vraag 3 na, één enkele groep --- toon
        aan dat dit $D_6$ is, bijvoorbeeld door een element van
        orde $6$ en een spiegelingsachtige involutie aan te wijzen.
  \item Toon aan dat $\Z/12\Z$, $\Z/6\Z \times \Z/2\Z$, $D_6$,
        $A_4$ en $\mathrm{Dic}_3$ paarsgewijs niet isomorf zijn
        \emph{(tel de elementen van orde $2$, of gebruik $n_3$)}.
        Hiermee is orde $12$ afgehandeld.
\end{enumerate}

\medskip
\noindent\textbf{Deel V --- Synthese.}
\begin{enumerate}[resume]
  \item Stel de classificatietabel samen: voor elke orde $n \leq
        15$ de volledige lijst groepen op isomorfie na, met de
        aantallen $1, 1, 1, 2, 1, 2, 1, 5, 2, 2, 1, 5, 1, 2, 1$.
\end{enumerate}

\medskip
\noindent\textbf{Deel VI --- Verder: de groepen van orde $p^3$ met
$p$ oneven.} De analyse van orde $8$ uit Deel III heeft een fraai
analogon voor oneven priemgetallen, met één werkelijk nieuw
verschijnsel. Zij $p$ oneven priem en $G$ niet-abels van orde
$p^3$.
\begin{enumerate}[resume]
  \item Toon aan dat $\abs{Z(G)} = p$, dat $G/Z(G) \cong
        (\Z/p\Z)^2$ \emph{(een cyclisch quotiënt naar het \omterm{ex:b3:groups:actions}{centrum}
        dwingt abelsheid af: \cref{exo:b3:groups:2})}, en dat
        $D(G) = Z(G)$ \emph{(voor $D(G) \subseteq Z(G)$: gebruik
        dat $G/Z(G)$ abels is; voor de gelijkheid: $G$ is
        niet-abels, dus $D(G) \neq \{e\}$)}. Leid af dat elke
        \omterm{def:b3:groups:derived}{commutator} $[x, y] = xyx^{-1}y^{-1}$ centraal is en een
        orde heeft die $p$ deelt.
  \item (De sleutelidentiteit) Zij $x, y \in G$ en $z = [y, x]$,
        centraal. Bewijs met inductie naar $k$:
        \[
        (xy)^k = x^k y^k z^{k(k-1)/2} .
        \]
        \emph{(Schuif elke $y$ langs elke $x$; elke kruising kost
        één centrale factor $z$.)}
  \item Leid af dat voor oneven $p$ de afbeelding $\theta\colon x
        \mapsto x^p$ een groepsmorfisme van $G$ naar $Z(G)$ is
        \emph{(waarom is $x^p$ centraal? en waarom heeft
        $z^{p(p-1)/2} = e$ nodig dat $p$ oneven is?)}, en besluit
        dat $G$ exponent $p$ of $p^2$ heeft, waarbij de twee
        gevallen zich onderscheiden doordat $\theta$ al dan niet
        triviaal is.
  \item (Exponent $p$) Stel dat elk element aan $x^p = e$ voldoet.
        Kies $x, y$ waarvan de klassen $G/Z(G)$ voortbrengen en
        stel $z = [y, x]$. Toon aan dat $z \neq e$, dat elk element
        van $G$ zich op precies één manier als $x^ay^bz^c$ schrijft
        ($0 \leq a, b, c < p$), en dat de vermenigvuldiging
        volledig vastligt door de betrekkingen $x^p = y^p = z^p =
        e$, $z$ centraal en $[y, x] = z$. Ga na dat de
        \emph{Heisenberggroep}
        \[
        H_p = \left\{
        \begin{pmatrix} 1 & a & c\\ 0 & 1 & b\\ 0 & 0 & 1
        \end{pmatrix} : a, b, c \in \mathbb F_p \right\}
        \subseteq GL_3(\mathbb F_p)
        \]
        deze betrekkingen realiseert en exponent $p$ heeft (bereken
        $(I + N)^p$ voor strikt bovendriehoekse $N$, met $N^3 = 0$
        en $p \geq 3$): elke niet-abelse groep van orde $p^3$ en
        exponent $p$ is isomorf met $H_p$.
  \item (Exponent $p^2$) Stel dat een $r \in G$ orde $p^2$ heeft en
        stel $N = \langle r\rangle$, \omterm{def:b3:groups:normal}{normaal} (index $p$:
        \cref{exo:b3:groups:10}). Toon aan dat er een $s \notin N$
        bestaat met $s^p = e$ \emph{(neem een willekeurige $t
        \notin N$ en corrigeer die met vraag 20: $\theta(t) = t^p
        \in Z(G) \subseteq N$ --- verantwoord $Z(G) = \langle
        r^p\rangle$ --- en kies $a$ zodat $s = tr^{a}$ voldoet aan
        $s^p = e$; waar wordt gebruikt dat $p$ oneven is?)}. Toon
        aan dat $srs^{-1} = r^{1+p}$, eventueel na vervanging van
        $s$ door een macht, en besluit: er is precies \emph{één}
        niet-abelse groep van orde $p^3$ en exponent $p^2$, namelijk
        $\Z/p^2\Z \rtimes_\varphi \Z/p\Z$ met $\varphi(1)\colon r
        \mapsto r^{1+p}$ \emph{(gebruik vraag 3;
        $\operatorname{Aut}(\Z/p^2\Z)$ is cyclisch van orde
        $p(p-1)$, hier zonder bewijs aangenomen, en heeft dus een
        unieke deelgroep van orde $p$)}.
  \item Besluit de telling: voor oneven $p$ zijn er precies $5$
        groepen van orde $p^3$ (drie abelse, twee niet-abelse),
        net als voor $p = 2$ --- maar de twee niet-abelse zijn niet
        langer $D_4$ en $Q_8$. Wijs precies aan waar het
        argument voor oneven $p$ breekt bij $p = 2$: in de
        identiteit van vraag 19 is $z^{k(k-1)/2}$ voor $k = p = 2$
        gelijk aan $z^1 \neq e$, dus is kwadrateren geen morfisme
        --- en inderdaad heeft $Q_8$ één enkel element van orde
        $2$, terwijl $D_4$, van exponent $4$, er vijf heeft.
\end{enumerate}

\medskip
\noindent\textbf{Deel VII --- Aanvullingen.}
\begin{enumerate}[resume]
  \item Tel voor oneven $p$ de elementen van orde $p$ in elk van de
        twee niet-abelse groepen van orde $p^3$: toon aan dat $H_p$
        er precies $p^3 - 1$ heeft en $M_p = \Z/p^2\Z \rtimes
        \Z/p\Z$ precies $p^2 - 1$ \emph{(gebruik het morfisme
        $\theta$ uit vraag 20: bepaal eerst zijn beeld, dan de orde
        van zijn kern)}. Reken het geval $p = 3$ na: $26$ tegen
        $8$. Verklaar waarom geen enkel argument met een
        kwadrateringsmorfisme van dit type $D_4$ van $Q_8$ kan
        onderscheiden, en welke telling dat wel doet.
  \item Noem een geheel getal $n \geq 1$ \emph{cyclisch} als elke
        groep van orde $n$ cyclisch is. Toon aan dat $n$ niet
        cyclisch is zodra $p^2 \mid n$ voor zeker priemgetal $p$,
        of zodra $n$ priemdelers $p < q$ heeft met $p \mid q - 1$
        \emph{(geef in elk geval een niet-cyclische groep van orde
        $n$; gebruik voor het tweede geval Deel II)}. Leid af dat
        $n$ cyclisch afdwingt dat $\gcd(n, \varphi(n)) = 1$, waarbij
        $\varphi$ de indicator van Euler is, en controleer dat aan
        de tabel van vraag 17: onder $n \leq 15$ zijn de ordes met
        één enkele groep precies $n \in \{1, 2, 3, 5, 7, 11, 13,
        15\}$, en dat zijn precies de $n$ met $\gcd(n, \varphi(n))
        = 1$.
\end{enumerate}
\end{problem}
